\section{Pattern and Association Rule Mining}
\subsection{Association Rule Mining}

\textbf{Association rule mining} aims to identify relationships between variables in large datasets. An association rule has the form \(X \rightarrow Y\), where \(X\) (antecedent) and \(Y\) (consequent) are itemsets. An itemset is a collection of one or more items, and a \(k\)-itemset contains exactly \(k\) items. Two key metrics evaluate association rules:

\begin{itemize}
    \item \textbf{Support}: The support of an itemset \(I\) is the fraction of transactions that contain \(I\):
    \[
    \text{support}(I) = \frac{\text{transactions containing }I}{\text{total transactions}}
    \]
    The \textbf{support count} \(\sigma(I)\) is the frequency of occurrence of itemset \(I\).
    
    \item \textbf{Confidence}: For an association rule \(X \rightarrow Y\), confidence measures how often \(Y\) appears in transactions that contain \(X\):
    \[
    \text{confidence}(X \rightarrow Y) = \frac{\text{support}(X \cup Y)}{\text{support}(X)}
    \]
\end{itemize}

A \textbf{frequent itemset} is an itemset whose support exceeds a predefined threshold, known as the \textbf{minimum support} (minsup).

\paragraph{Association Rule Mining Strategies}: The goal is to find rules whose support and confidence exceed the minsup and minconf thresholds, respectively. A brute-force approach computes support and confidence for every possible rule, then prunes those that miss the thresholds. A more efficient two-step approach is:
\begin{itemize}
    \item \textbf{Frequent Itemset Generation}: Identify itemsets with support \(\geq\) minsup.
    \item \textbf{Rule Generation}: Create high-confidence rules from each frequent itemset.
\end{itemize}

Despite this decomposition, frequent itemset generation remains costly, prompting the development of efficient algorithms like Apriori and FP-Growth.

\subsection{The Apriori Algorithm}

The Apriori algorithm is a key method for mining frequent itemsets and generating association rules. It relies on the Apriori principle: if an itemset is frequent, then all of its subsets must also be frequent. This follows from the anti-monotone property of support: the support of an itemset never exceeds the support of its subsets.

\begin{algorithm}
\caption{Apriori Algorithm}
\begin{algorithmic}[1]
\State Let $k=1$
\State Generate $F_1 =$ \{frequent 1-itemsets\}
\Repeat
    \State \textbf{Candidate Generation:} Generate $L_{k+1}$ from $F_k$
    \State \textbf{Candidate Pruning:} Prune candidate itemsets in $L_{k+1}$ containing subsets of length $k$ that are infrequent
    \State \textbf{Support Counting:} Count the support of each candidate in $L_{k+1}$ by scanning the database
    \State \textbf{Candidate Elimination:} Eliminate candidates in $L_{k+1}$ that are infrequent, leaving only those that are frequent $\Rightarrow F_{k+1}$
    \State $k = k + 1$
\Until{$F_k$ is empty}
\end{algorithmic}
\end{algorithm}

The Apriori algorithm iteratively identifies frequent itemsets of increasing size. At each iteration, it generates candidate itemsets of size $k+1$ from frequent itemsets of size $k$.

\begin{itemize}
    \item \textbf{Candidate Generation}: The method merges frequent $(k-1)$-itemsets to generate candidate itemsets:
    \begin{itemize}
        \item For $k=2$: merge any two distinct frequent $1$-itemsets (the empty $(k-2)$-prefix condition is trivially met), producing the candidate $2$-itemset $\{a,b\}$.
        \item For $k > 2$: merge two frequent $(k-1)$-itemsets whose first $(k-2)$ items are identical.
    \end{itemize}

    \item \textbf{Candidate Pruning}: Prune candidates containing any subset that is not frequent.

    \item \textbf{Support Counting and Candidate Elimination}: Count the support of each candidate itemset by scanning the database. Remove candidates with support below the minsup threshold, leaving only frequent itemsets.

    \item \textbf{Rule Generation}: For a frequent itemset $L$, generate association rules from its non-empty subsets $f \subset L$. A rule $f \rightarrow L - f$ is valid if it meets the minimum confidence requirement. For a frequent itemset of size $k$, there are $2^k - 2$ candidate association rules.
\end{itemize}

\subsection{Pattern Evaluation and Interestingness Measures}

Association rule algorithms often generate many rules, not all of which are useful. Interestingness measures can help prune or rank patterns. These measures are computed using a contingency table:

\begin{table}[h!]
\centering
\begin{tabular}{|c|c|c|c|}
\hline
 & $Y$ & $\neg Y$ & \\
\hline
$X$ & $f_{11}$ & $f_{10}$ & $f_{1+}$ \\
\hline
$\neg X$ & $f_{01}$ & $f_{00}$ & $f_{0+}$ \\
\hline
 & $f_{+1}$ & $f_{+0}$ & $N$ \\
\hline
\end{tabular}
\caption{Contingency Table for Association Rules}
\end{table}

Where:
\begin{itemize}
    \item $f_{11}$: Support of $X$ and $Y$ (transactions containing both $X$ and $Y$)
    \item $f_{10}$: Support of $X$ and $\neg Y$ (transactions containing $X$ but not $Y$)
    \item $f_{01}$: Support of $\neg X$ and $Y$ (transactions containing $Y$ but not $X$)
    \item $f_{00}$: Support of $\neg X$ and $\neg Y$ (transactions containing neither $X$ nor $Y$)
    \item $N$: Total number of transactions
\end{itemize}

\textbf{Statistical relationship between items}: we compare the confidence with the support of the consequent:
\begin{itemize}
    \item If $\text{confidence}(X \rightarrow Y) > \text{support}(Y)$, $X$ and $Y$ are positively correlated.
    \item If $\text{confidence}(X \rightarrow Y) = \text{support}(Y)$, $X$ and $Y$ are independent.
    \item If $\text{confidence}(X \rightarrow Y) < \text{support}(Y)$, $X$ and $Y$ are negatively correlated.
\end{itemize}

The same relationship holds in terms of joint and marginal probabilities:
\begin{itemize}
    \item $P(X,Y) > P(X) \times P(Y)$ implies $X$ and $Y$ are positively correlated.
    \item $P(X,Y) = P(X) \times P(Y)$ implies $X$ and $Y$ are independent.
    \item $P(X,Y) < P(X) \times P(Y)$ implies $X$ and $Y$ are negatively correlated.
\end{itemize}

\begin{definition}
\textbf{Lift} (for rules) and \textbf{Interest} (for itemsets) are measures that account for statistical dependence:
\[
\text{lift}(X \rightarrow Y) = \frac{\text{confidence}(X \rightarrow Y)}{\text{support}(Y)} = \frac{P(X,Y)}{P(X) \times P(Y)}
\]
A lift value greater than 1 indicates positive correlation, exactly 1 indicates independence, and less than 1 indicates negative correlation.
\end{definition}

\subsubsection{Handling Continuous and Categorical Attributes}
Association rule mining traditionally targets binary attributes, but it extends to categorical and continuous attributes as well.

\paragraph{Categorical Attributes} A common approach is to create a new ``item'' for each distinct attribute-value pair. When attributes have many values, those with low support might be aggregated.

\paragraph{Continuous Attributes} Methods for handling continuous attributes include:
\begin{itemize}
    \item Discretization-based methods (convert continuous values into intervals),
    \item Statistics-based methods (use statistical measures),
    \item Non-discretization methods like minApriori.
\end{itemize}

\paragraph{Multi-level Association Rules} Rules can be mined at several levels of an item taxonomy:
\begin{itemize}
    \item rules at lower levels might lack support and be too specific,
    \item higher-level rules give a more general view of the relationships.
\end{itemize}
Two approaches handle these rules:
\begin{itemize}
    \item augment each transaction with its higher-level items, then mine as usual,
    \item generate frequent patterns starting at the highest level of the taxonomy.
\end{itemize}
The first increases data dimensionality, while the second increases I/O and could miss cross-level patterns.

\subsubsection{Performance Issues with Apriori}

The Apriori algorithm has performance limitations:
\begin{itemize}
    \item Candidate generation can lead to large sets, e.g., $10^4$ frequent 1-itemsets generating more than $10^7$ candidate 2-itemsets.
    \item Discovering a pattern of size 100 requires generating $2^{100} \approx 10^{30}$ candidates.
    \item Apriori requires $(n+1)$ database scans, where $n$ is the pattern length.
\end{itemize}
These challenges led to more efficient algorithms, such as FP-Growth.


\subsection{FP-Growth Algorithm}

\textbf{FP-Growth (Frequent Pattern Growth)} mines frequent patterns without generating candidates. It compresses the database into a compact structure, the Frequent-Pattern tree (FP-tree), and splits the mining task into smaller subtasks.

\begin{algorithm}
\caption{FP-tree Construction}
\begin{algorithmic}[1]
\Procedure{BuildFPTree}{$Dataset~D, minsup$}
    \State Scan $D$ to find frequent 1-itemsets and their supports
    \State Sort frequent items by support to create list $L$
    \State Create the root of FP-tree labeled "null"
    \For{each transaction $t$ in $D$}
        \State Filter $t$ to keep only frequent items, sort by $L$
        \State Insert the filtered transaction into the tree
    \EndFor
    \State \Return FP-tree
\EndProcedure
\end{algorithmic}
\end{algorithm}

FP-tree construction guarantees \textbf{completeness} (preserves full transaction patterns) and \textbf{compactness} (removes infrequent items, orders by frequency).

\begin{algorithm}
\caption{FP-Growth Algorithm}
\begin{algorithmic}[1]
\Procedure{FPGrowth}{$FP\text{-}tree, \alpha$}
    \For{each frequent item $a_i$ in the header table}
        \State Construct conditional pattern base for item $a_i$
        \State Build conditional FP-tree $FP\text{-}tree_{a_i}$
        \If{$FP\text{-}tree_{a_i} \neq \emptyset$} \Comment{Recursively mine conditional FP-trees}
            \If{$FP\text{-}tree_{a_i}$ has a single path} \Comment{enumerate all patterns}
                \State Generate all combinations of items in the path with $\alpha$
            \Else
                \State Call \Call{FPGrowth}{$FP\text{-}tree_{a_i}, a_i \cup \alpha$}
            \EndIf
        \EndIf
    \EndFor
\EndProcedure
\end{algorithmic}
\end{algorithm}

The algorithm constructs a conditional pattern base by following node-links and accumulating prefix paths for each item.

Two properties of the FP-tree make this efficient:
\begin{itemize}
    \item \textbf{Node-link property:} frequent patterns can be obtained by following node-links.
    \item \textbf{Prefix path property:} only the prefix paths of a node need to be accumulated, each with the same frequency count.
\end{itemize}

For each pattern base, we accumulate the item counts and build a conditional FP-tree. When that conditional tree reduces to a single path, we generate the frequent patterns by enumerating all combinations of items along the path.

\paragraph{Advantages of FP-Growth}

FP-Growth outperforms Apriori for several reasons:
\begin{itemize}
    \item No candidate generation.
    \item Compact FP-tree structure.
    \item Fewer database scans.
    \item Simpler operations (counting and tree building).
\end{itemize}

\subsubsection{Special Types of Frequent Patterns}

In addition to regular frequent itemsets, several special types of frequent patterns may be of interest:

\begin{definition}
An itemset is \textbf{maximal frequent} if it is frequent and none of its immediate supersets are frequent. Maximal frequent itemsets form a boundary between frequent and infrequent itemsets in the itemset lattice.
\end{definition}

\begin{definition}
An itemset $X$ is \textbf{closed} if no immediate superset of $X$ has the same support as $X$. $X$ is not closed if at least one of its immediate supersets has the same support count as $X$.
\end{definition}

Maximal frequent itemsets are a subset of closed frequent itemsets, which are, in turn, a subset of all frequent itemsets. Mining maximal or closed frequent itemsets can significantly reduce the number of patterns considered while preserving essential information.
